\chapter{УДИРТГАЛ}

\section{Үндэслэл, ач холбогдол}

Хөдөө орон нутагт рашаан, шавар эмчилгээний үйл ажиллагаа явуулдаг 20 гаруй сувиллууд байдаг. Тэдгээр сувиллуудын өнөөгийн захиалга авах үйл явц дараах алхмуудаас бүрдэж байна. Үүнд:

\begin{itemize}
    \item Үйлчлүүлэгч тухайн сувиллын хүлээн авах утас руу залгана.
    \item Хүлээн авах ажилтан захиалгын мэдээллийг цаасан дээр гараар тэмдэглэнэ.
    \item Хэвтэн эмчлүүлэх захиалга товлосон өдөр дөхөх үед, ажилтан үйлчлүүлэгч рүү эргэж залган урьдчилгаа төлбөр авах, ирэх эсэхийг дахин баталгаажуулна.
\end{itemize}

Энэ урсгал нь удаан хугацаанд хэрэглэгдэж ирсэн ч дараах асуудлуудыг үүсгэж байна:

\begin{itemize}
    \item Үйлчлүүлэгч захиалгаа урьдчилан мэдэгдэлгүй цуцлах, эсвэл хоорондын ойлголцлын зөрөөнөөс шалтгаалж огноо, хүмүүсийн тоо зэрэг алдаа гарах.
    \item Хүлээн авах ажилтны гар тэмдэглэл алдагдах, гээх зэрэг хүний хүчин зүйлээс шалтгаалсан эрсдэл.
    \item Утасны ачааллаас шалтгаалан шуурхай хариу өгөх боломжгүй байх.
    \item Тайлан тооцоог нарийвчлалтай гаргахад хүндрэлтэй.
\end{itemize}

Иймээс хэзээ ч, хаанаас ч вэбээр захиалга өгөх боломжтой амрагчийн захиалгын системийг нэвтрүүлэх хэрэгцээ бодитойгоор бий болсон гэж үзсэн. Энэхүү судалгааны ажлын үр дүн нь тухайн нэг байгууллага төдийгүй ижил төрлийн бусад сувиллын байгууллагуудад ч хэвшүүлэх, ашиглах боломжтой жишиг загвар болох ач холбогдолтой билээ.

\section{Судалгааны ажлын зорилго}

Энэхүү судалгааны ажлын ерөнхий зорилго нь Өвөрхангай аймгийн Хужирт сумын сувиллын онцлогт тохируулсан амрагчийн захиалгын бүртгэл болон захиалгыг хянах админ системийг вэб хэлбэрээр боловсруулж, захиалга авах үйл явцыг илүү үр ашигтай, алдаа багатай, ашиглахад хялбар болгох юм.

\section{Судалгааны ажлын зорилтууд}

Дээрх зорилгыг хэрэгжүүлэхийн тулд дараах зорилтуудыг дэвшүүлж байна.

\begin{enumerate}
    \item Хүлээн авах ажилтан болон удирдлагад тулгарч буй асуудал, одоогийн захиалгын урсгал, түүнээс үүдэлтэй хохирол, эрсдэлийг тодорхойлох.
    \item Сувиллын үйлчлүүлэгч, ажилтнуудын хэрэгцээнд нийцсэн вэб суурьтай захиалгын системийн шаардлага, архитектур, интерфейсийг тодорхойлох.
    \item Захиалга авах, баталгаажуулах, админаар хянах, тайлан, статистик гаргах боломж бүхий вэб системийн frontend болон backend хэсгийг хөгжүүлж, PostgreSQL өгөгдлийн сантай уялдуулан туршилтын хувилбарыг програмчлах.
    \item Хөгжүүлсэн системийг cloud server (AWS) орчинд байршуулж, жинхэнэ эсвэл туршилтын хэрэглэгчдээр ашиглуулан туршилт, үнэлгээ хийж, санал хүсэлтийн дагуу сайжруулалт хийх.
\end{enumerate}
\section{Судалгааны асуулт ба хэмжигдэх шалгуур}

Судалгааны ажлын зорилтуудыг хэмжихүйц шалгууруудтай холбохын тулд дараах судалгааны асуултуудыг тодорхойлов. Асуулт бүрийн хариуг баруун баганад буй тодорхой хэмжүүрээр үнэлнэ.

\begin{table}[H]
\centering
\caption{Судалгааны асуулт ба үнэлгээний шалгуур үзүүлэлтүүд}
\label{tab:research-questions}
\begin{tabular}{|c|p{5.5cm}|p{6cm}|}
\hline
\textbf{RQ} & \textbf{Судалгааны асуулт} & \textbf{Хэмжих шалгуур үзүүлэлт} \\
\hline
RQ1 & Вэб суурьтай захиалгын систем нь цаасан бүртгэл, утсаар захиалга авах процесстой харьцуулахад захиалга боловсруулах хугацааг бууруулж чадах уу? & Нэг захиалга боловсруулах дундаж хугацаа (секунд), мэдээлэл хайх хугацаа, тайлан гаргах хугацааны харьцуулалт \\
\hline
RQ2 & Систем нь давхар захиалга (double-booking), огнооны зөрөө зэрэг хүний алдаанаас үүдэлтэй асуудлуудыг арилгаж чадах уу? & Зэрэг хандалтын туршилтаар илэрсэн давхар захиалгын тоо, pessimistic locking-ийн амжилтын хувь \\
\hline
RQ3 & Системийн ашиглах чадвар (usability) хүлээн зөвшөөрөгдөхүйц түвшинд байх уу? & System Usability Scale (SUS) оноо 68-аас дээш байх, даалгавар гүйцэтгэх амжилтын хувь \\
\hline
\end{tabular}
\end{table}